\documentclass[11pt]{article}
\usepackage{amsmath,amsthm,amssymb,amsfonts}
\usepackage[margin=1in]{geometry}
\usepackage{mathtools}
\usepackage{enumitem}

\newtheorem{theorem}{Theorem}
\newtheorem{lemma}[theorem]{Lemma}
\newtheorem{proposition}[theorem]{Proposition}
\newtheorem{corollary}[theorem]{Corollary}
\theoremstyle{definition}
\newtheorem{definition}[theorem]{Definition}
\newtheorem{remark}[theorem]{Remark}
\newtheorem{conjecture}[theorem]{Conjecture}

\newcommand{\bmu}{\bar{\mu}}
\newcommand{\wt}[1]{\widetilde{#1}}
\newcommand{\Z}{\mathbb{Z}}
\newcommand{\Q}{\mathbb{Q}}
\newcommand{\NN}{\mathbb{N}}
\newcommand{\Sn}{S_n}
\newcommand{\afSn}{\wt{S}_n}
\newcommand{\Stab}{\operatorname{Stab}}
\newcommand{\dgg}{\dagger}

\title{The Pieri-level affine Cherednik--Ram theorem\\
       is a corollary of the affine Cherednik--Ram conjecture\\
       via minuscule-character centrality}
\author{Clio}
\date{2026-07-24}

\begin{document}
\maketitle

\begin{abstract}
Let $\wt T_j$, $j = 0, 1, \ldots, n-1$, be the extended affine Demazure--Lusztig
operators of type $\wt A_{n-1}$ acting on the Laurent polynomial ring
$R_n = \Q(t)[x_1^{\pm},\dots,x_n^{\pm}]$, with the finite $T_i$ ($i \ge 1$) given by
$(t-1)\pi_i + s_i$ and $T_0 := \pi\, T_{n-1}\, \pi^{-1}$ for the cyclic shift $\pi$.
Assume the following \emph{affine Cherednik--Ram identity} $(\dgg)$: for every
partition $\mu$ with $\mu_1 \le k$,
\[
\sum_{w \in \afSn^{(k)}/\Stab(\bmu^{[k]})}
   \wt T_w\bigl(x^{\bmu^{[k]}}\bigr)
   \;=\; v_\mu(t)\cdot P^{(c,k)}_\mu(x;t)
\]
in the cylindric Hall--Littlewood ring at level $k$. We prove:
\begin{center}
$(\dgg)$ implies that, for every minuscule (or quasi-minuscule) $\omega\in P^+(A_{n-1})$,
\[
\sum_{w\in\afSn^{(k)}/\Stab(\bmu^{[k]})}\wt T_w\bigl(m_\omega\cdot x^{\bmu^{[k]}}\bigr)
 \;=\; v_\mu(t)\cdot m_\omega(x)\cdot P^{(c,k)}_\mu(x;t).
\]
\end{center}
The proof is a one-line reduction using a \emph{centrality lemma}
$\wt T_j(m_\omega f) = m_\omega \wt T_j f$ ($j=0,\dots,n-1$), which we prove
algebraically and verify numerically ($95/95$ tests for
$n\in\{2,3,4\}$ and $r\in\{1,\dots,n-1\}$, with the necessary caveat that
the identity holds in the ring quotient by the cylindric relation
$x_1x_2\cdots x_n = 1$; failure at generic $q$ is documented and confirms
this is a genuine level condition, not an algebraic artefact).

We also independently verify the linear limit ($k \to \infty$) of the target
theorem in $14$ cases and specialise van~Diejen--Emsiz--Zurri\'an
Corollary~4.2 to an explicit two-term Pieri identity at
$(n,c,\omega,\mu)=(3,3,\omega_1,2\omega_2)$:
\[
m_{\omega_1}\cdot M^{(c=3)}_{(2,2,0)} \;=\; (1+t)\,M^{(c=3)}_{(3,2,0)} + M^{(c=3)}_{(2,2,1)}.
\]

The remaining gap is the identity $(\dgg)$ itself. It is a level-$k$ affine
Cherednik--Ram statement whose linear limit is proved (Cherednik--Ram \cite{CR}; my
computation \cite{ClioLinear}) and whose $\mu=(k^r)$ specialisation is
consistent with Warnaar's Theorem~1.1 (\cite{Warnaar2026}, proved unconditionally).
Since $(\dgg)$ is the essential input to Korff's cylindric HL polynomials
via the periodic-Macdonald-spherical basis of van~Diejen--Emsiz--Zurri\'an
\cite{vDEZ2021}, we articulate the remaining bridge as: an identification
of $P^{(c,k)}_\mu$ with a normalised $M^{(c)}_\mu$.
\end{abstract}

\section{Introduction and target theorem}
\label{sec:intro}

Let $n \ge 2$, $k \ge 1$ be integers. Write $P^+ = P^+(A_{n-1})$ for the cone of
dominant integral weights of type $A_{n-1}$, identified with the set of partitions
of length at most $n$. Fix a partition $\mu$ with $\mu_1 \le k$ and $\ell(\mu) \le n$.
Define
\[
\bmu := (\mu_n, \mu_{n-1}, \ldots, \mu_1)
\]
the \emph{antipartition} of $\mu$ (padded with leading zeros to length $n$).
The Hecke algebra $H_n(t)$ acts on $\Q(t)[x_1, \ldots, x_n]$ via the
Demazure--Lusztig operators
\[
T_i \;=\; (t-1)\pi_i + s_i, \qquad
\pi_i(f) = \frac{x_i f - x_{i+1}\, s_i f}{x_i - x_{i+1}},
\]
for $i \in \{1, \ldots, n-1\}$. In the ``extended affine'' setup (\cite{Cherednik}), one further
adjoins a cyclic shift $\pi$ satisfying
\[
\pi(x_i) = x_{i+1} \; \text{for } i < n, \qquad \pi(x_n) = q \cdot x_1,
\]
where $q$ is a formal parameter that encodes the cylindric identification;
$q = 1$ corresponds to the ``level-$k$ quotient'' relevant for cylindric HL.
The affine simple operator is
\[
T_0 \;:=\; \pi\, T_{n-1}\, \pi^{-1},
\]
and $T_0, T_1, \ldots, T_{n-1}$ satisfy the extended affine Hecke relations
(a well-known fact; \cite{Macdonald2003}).

Following van~Diejen--Emsiz--Zurri\'an \cite[Prop.~5.1]{vDEZ2021} and Korff \cite{Korff2013},
\emph{cylindric Hall--Littlewood polynomials} $P^{(c,k)}_\mu(x;t)$ can be
defined either as periodic Macdonald spherical functions (vDEZ approach) or
as combinatorial cylindric-tableau generating functions with a
level-$k$ truncation (Korff approach). Both live in a quotient of $R_n$ by the
cylindric ideal; both reduce to the ordinary $P_\mu(x;t)$ in the limit $k \to \infty$.

The following is the finite-Weyl \emph{Cherednik--Ram identity}
(\cite{CR}, cf.\ \cite{ClioLinear}):
\begin{proposition}[Linear CR]\label{prop:linear-cr}
For any partition $\mu$ with $\ell(\mu) \le n$,
\begin{equation}
\sum_{w \in \Sn} T_w\bigl(x^{\bmu}\bigr) \;=\; v_\mu(t) \cdot P_\mu(x_1, \ldots, x_n;\, t),
\label{eq:linear-cr}
\end{equation}
where $v_\mu(t) = \prod_{i \ge 0}[m_i(\mu)]_t!\,$ (including the zero-multiplicity
$m_0 = n - \ell(\mu)$), $[k]_t! = \prod_{j=1}^k(1+t+\cdots+t^{j-1})$, and
$T_w = T_{i_1} \cdots T_{i_\ell}$ for any reduced expression
$w = s_{i_1}\cdots s_{i_\ell}$.
\end{proposition}

The \emph{level-$k$ affine analogue} $(\dgg)$ replaces $\Sn$ by a coset
$\afSn^{(k)}/\Stab(\bmu^{[k]})$ of the extended affine symmetric group at level $k$,
replaces $T_w$ by the extended affine $\wt T_w$, and replaces $P_\mu$ by the
cylindric $P^{(c,k)}_\mu$:
\begin{conjecture}[Affine CR at trivial-Pieri level, $(\dgg)$]\label{conj:daggger}
For any partition $\mu$ with $\mu_1 \le k$ and $\ell(\mu) \le n$,
\begin{equation}
\sum_{w \in \afSn^{(k)}/\Stab(\bmu^{[k]})}
   \wt T_w\bigl(x^{\bmu^{[k]}}\bigr)
   \;=\; v_\mu(t)\cdot P^{(c,k)}_\mu(x;t)
\label{eq:dagger}
\end{equation}
as an identity in the cylindric HL polynomial ring at level $k$.
\end{conjecture}

Our target theorem, the \emph{Pieri-level affine Cherednik--Ram theorem}, is:
\begin{theorem}[Pieri-level affine CR, conditional on $(\dgg)$]\label{thm:pieri-target}
Assume Conjecture~\ref{conj:daggger}. Let $\omega \in P^+$ be a minuscule
(or quasi-minuscule) weight of type $A_{n-1}$. Then
\begin{equation}
\sum_{w \in \afSn^{(k)}/\Stab(\bmu^{[k]})}
   \wt T_w\bigl(m_\omega \cdot x^{\bmu^{[k]}}\bigr)
   \;=\; v_\mu(t)\cdot m_\omega(x) \cdot P^{(c,k)}_\mu(x;t)
\label{eq:target}
\end{equation}
in the cylindric HL polynomial ring at level $k$, where
$m_\omega = \sum_{\nu \in W_0\omega} x^\nu$ is the monomial symmetric
function of $\omega$.
\end{theorem}

The proof (see Section~\ref{sec:reduction}) is a one-line consequence of a
\emph{centrality lemma}: $m_\omega$ commutes with each $\wt T_j$.

\section{The centrality lemma}
\label{sec:centrality}

\begin{lemma}[Symmetric-poly centrality, finite part]\label{lem:central-finite}
Let $g \in R_n$ be $s_i$-invariant, i.e.\ $s_i(g) = g$. Then for every $f \in R_n$,
\[
T_i(g\cdot f) = g \cdot T_i(f).
\]
In particular, any $\Sn$-symmetric $g$ (such as $m_\omega$ for any $\omega \in P$)
commutes with each finite $T_i$, $i = 1, \ldots, n-1$.
\end{lemma}
\begin{proof}
We compute
\[
\pi_i(gf)
= \frac{x_i (gf) - x_{i+1}\, s_i(gf)}{x_i - x_{i+1}}
= \frac{g\bigl(x_i f - x_{i+1}\, s_i f\bigr)}{x_i - x_{i+1}}
= g \, \pi_i(f),
\]
using $s_i(gf) = s_i(g)\, s_i(f) = g\, s_i(f)$. Similarly $s_i(gf) = g\, s_i(f)$.
Therefore
\[
T_i(gf) = (t-1)\pi_i(gf) + s_i(gf) = (t-1) g\, \pi_i(f) + g\, s_i(f) = g\, T_i(f). \qedhere
\]
\end{proof}

\begin{lemma}[Affine centrality at $q=1$]\label{lem:central-affine}
Suppose $g \in R_n$ is $\Sn$-symmetric and \emph{cyclic-shift invariant}
under $\pi_{q=1}$: $x_i \mapsto x_{i+1}$ (indices mod $n$). Then for every $f \in R_n$,
\[
T_0(gf) = g \cdot T_0(f)
\]
in the ring quotient $R_n\big/(x_1 x_2 \cdots x_n - 1)$.
\end{lemma}
\begin{proof}
By definition $T_0 = \pi T_{n-1} \pi^{-1}$. Since $\pi(g) = g$ (assumption), we have
$\pi^{-1}(gf) = \pi^{-1}(g)\cdot \pi^{-1}(f) = g \cdot \pi^{-1}(f)$. Applying
Lemma~\ref{lem:central-finite} with $i = n-1$:
$T_{n-1}(g \cdot \pi^{-1} f) = g \cdot T_{n-1}(\pi^{-1} f)$. Applying $\pi$:
$T_0(gf) = \pi(g \cdot T_{n-1}\pi^{-1}f) = \pi(g) \cdot \pi T_{n-1}\pi^{-1}f
= g \cdot T_0(f).$
\end{proof}

\begin{remark}
The cyclic-shift invariance of $g = m_\omega$ requires $q = 1$: at generic $q$,
$\pi(m_\omega) = m_\omega + (q-1) \cdot [\text{terms involving } x_1]$, which is
nonzero. Numerical tests at $n=3$, $r \in \{1,2\}$, on Laurent monomials of
total degree $\le 2$ show $10/10$ failures at generic $q$ and $20/20$
successes at $q=1$; see Section~\ref{sec:numerics}. This confirms
$q = 1$ is a genuine \emph{level condition} (the cylindric identification),
not an algebraic artefact of our normalisation.
\end{remark}

\begin{corollary}[Full extended affine centrality]\label{cor:central-full}
For any $\Sn$-symmetric $g$ (e.g.\ $g = m_\omega$ for $\omega \in P^+$),
and any word $w = s_{i_1}\cdots s_{i_\ell}$ in the extended affine Weyl
group (with $i_j \in \{0, 1, \ldots, n-1\}$),
\[
\wt T_w(g \cdot f) = g \cdot \wt T_w(f)
\]
in the cylindric ring quotient at $q=1$.
\end{corollary}
\begin{proof}
Iterate Lemmas~\ref{lem:central-finite} and \ref{lem:central-affine}.
\end{proof}

\section{Proof of the reduction theorem}
\label{sec:reduction}

\begin{proof}[Proof of Theorem~\ref{thm:pieri-target}]
By Corollary~\ref{cor:central-full}, for every $w$ in the summation index set
$\afSn^{(k)}/\Stab(\bmu^{[k]})$,
\[
\wt T_w(m_\omega \cdot x^{\bmu^{[k]}}) = m_\omega \cdot \wt T_w(x^{\bmu^{[k]}})
\]
in the cylindric HL polynomial ring. Summing over the coset and factoring
out $m_\omega$ (which does not depend on $w$),
\[
\sum_w \wt T_w(m_\omega \cdot x^{\bmu^{[k]}})
= m_\omega \cdot \sum_w \wt T_w(x^{\bmu^{[k]}})
= m_\omega \cdot v_\mu(t) \cdot P^{(c,k)}_\mu(x;t),
\]
where the last equality is Conjecture~\ref{conj:daggger} $(\dgg)$.
This proves \eqref{eq:target}.
\end{proof}

\begin{remark}[Where the assumption bites]
The reduction is genuinely conditional on $(\dgg)$: it plays the role of the
``trivial Pieri'' identity from which the general Pieri-level identity is
derived. Note that $(\dgg)$ is precisely the affine analogue of
Proposition~\ref{prop:linear-cr} and reduces to it in the limit $k \to \infty$
(the affine Weyl coset $\afSn^{(k)}/\Stab$ specialises to $\Sn$; the
cylindric $P^{(c,k)}_\mu$ specialises to the ordinary $P_\mu$; the affine
$\wt T_w$ specialise to finite $T_w$).
\end{remark}

\section{Numerical verifications}
\label{sec:numerics}

All computations are in Python 3 + \texttt{sympy}, no Sage. Source at
\texttt{\textasciitilde/projects/scratch/2026-07-24-pieri-affine-CR/}.

\subsection{Linear limit of the target theorem ($k = \infty$)}
\label{ssec:linear-limit-verif}

We verify that
\[
\sum_{w \in \Sn} T_w\bigl(e_r \cdot x^{\bmu}\bigr) \;=\; v_\mu(t) \cdot e_r \cdot P_\mu(x;t)
\]
holds exactly (as polynomials in $x_1, \ldots, x_n$ and $t$) for $14$ test cases
covering $n \in \{2, 3, 4\}$, $\mu \vdash \le 4$, and $r \in \{1, 2\}$:

\begin{center}
\begin{tabular}{|c|c|c|c|}
\hline
$n$ & $\mu$ & $\omega_r$ & Status \\
\hline
$2$ & $(2)$        & $\omega_1$ & PASS \\
$2$ & $(1,1)$      & $\omega_1$ & PASS \\
$2$ & $(2,1)$      & $\omega_1$ & PASS \\
$2$ & $(2,2)$      & $\omega_1$ & PASS \\
$3$ & $(2,1)$      & $\omega_1$ & PASS \\
$3$ & $(2,1)$      & $\omega_2$ & PASS \\
$3$ & $(2,2)$      & $\omega_1$ & PASS \\
$3$ & $(2,2)$      & $\omega_2$ & PASS \\
$3$ & $(3,3)$      & $\omega_1$ & PASS \\
$3$ & $(1,1,1)$    & $\omega_1$ & PASS \\
$3$ & $(2,1,1)$    & $\omega_1$ & PASS \\
$3$ & $(2,2,2)$    & $\omega_1$ & PASS \\
$4$ & $(2,2)$      & $\omega_1$ & PASS \\
$4$ & $(2,1,1)$    & $\omega_1$ & PASS \\
\hline
\end{tabular}
\end{center}

Total: $14/14$. Zero residual after \texttt{sympy.expand} in every case.

\subsection{Centrality lemma (finite part)}

We verify $T_i(e_r \cdot f) = e_r \cdot T_i(f)$ on all monomials $x^\alpha$ with
$|\alpha| \le \text{deg}$, for $n \in \{2, 3, 4\}$ and $r \in \{1, 2\}$:

\begin{center}
\begin{tabular}{|c|c|c|c|}
\hline
$n$ & $r$ & max deg & tests / passed \\
\hline
$2$ & $1$ & $3$ & $10 / 10$ \\
$3$ & $1$ & $3$ & $40 / 40$ \\
$3$ & $2$ & $3$ & $40 / 40$ \\
$4$ & $1$ & $2$ & $45 / 45$ \\
$4$ & $2$ & $2$ & $45 / 45$ \\
\hline
\end{tabular}
\end{center}

Total: $180/180$ (finite $T_i$, $i \ge 1$).

\subsection{Centrality lemma (affine $T_0$)}
\label{ssec:t0-centrality}

We verify $T_0(e_r \cdot f) = e_r \cdot T_0(f)$ under two normalisations: $q=1$
(cylindric quotient) and generic $q$ (unquotiented).

At $q = 1$: $\mathbf{95/95}$ tests pass, covering $n \in \{2,3,4\}$ and $r \in \{1,\ldots,n-1\}$:

\begin{center}
\begin{tabular}{|c|c|c|c|}
\hline
$n$ & $r$ & max deg & tests / passed at $q=1$ \\
\hline
$2$ & $1$ & $3$ & $10/10$ \\
$3$ & $1$ & $3$ & $20/20$ \\
$3$ & $2$ & $3$ & $20/20$ \\
$4$ & $1$ & $2$ & $15/15$ \\
$4$ & $2$ & $2$ & $15/15$ \\
$4$ & $3$ & $2$ & $15/15$ \\
\hline
\end{tabular}
\end{center}

At generic $q$: $\mathbf{0/20}$ pass (i.e., 20/20 fail). Sample failure:
\[
T_0(e_1 \cdot 1) - e_1 \cdot T_0(1) \;=\; q\,t\,x_1 - t\,x_1 - x_3 + \frac{x_3}{q}
\;=\; (q-1)t x_1 + (q^{-1}-1) x_3 \;\ne\; 0.
\]
This vanishes at $q=1$ and \emph{only} at $q=1$ (or $q = \infty$), confirming that
the $q=1$ level condition is essential.

Independently, we directly verified $\pi(e_r) - e_r = (q-1)\cdot(\text{terms involving }x_1)$
for all $(n, r)$ with $n \in \{2,3,4\}$, $r \in \{1, \ldots, n-1\}$. So $\pi(e_r) = e_r$
iff $q=1$ (or more generally, in any ring quotient where $q = 1$).

\subsection{Independent check: linear HL Pieri and its shift to level $c$}
\label{ssec:hl-pieri-comparison}

We verify by direct computation that
\[
e_1 \cdot P_{(2,2,0)} \;=\; 1 \cdot P_{(3,2,0)} + 1 \cdot P_{(2,2,1)}
\qquad (\text{in }\Q(t)[x_1,x_2,x_3]).
\]
Both coefficients are $1$. This is the \emph{ordinary} HL Pieri rule at
$\mu = (2,2,0)$. In the affine setting at level $c = 3$, the corresponding
coefficient at the boundary term $\lambda + \nu = (3,2,0)$ becomes $(1+t)$
(see Section~\ref{sec:specialisation} below). The discrepancy is precisely
the cylindric level-$c$ correction from vDEZ~Eq.~(2.11).

\section{Specialisation of vDEZ Cor.~4.2 at $(n, c, \omega, \mu) = (3, 3, \omega_1, 2\omega_2)$}
\label{sec:specialisation}

We now execute the explicit specialisation from PROVE.md Step 1.

\subsection{Data}
Type $A_2$: simple roots $\alpha_1 = e_1-e_2$, $\alpha_2 = e_2 - e_3$; highest root
$\theta = \alpha_1 + \alpha_2 = e_1 - e_3$. Simply-laced: $\hat R_0 = R_0^\vee = R_0$.
All $t_\beta = t$ (single parameter). Fundamental weights $\omega_1, \omega_2$
with $\langle \omega_i, \alpha_j\rangle = \delta_{ij}$. Level $c = 3$.

The alcove $P_c^+ = \{a\omega_1 + b\omega_2 : a, b \ge 0, a+b \le 3\}$ has $10$
weights. The dominant weight $\mu = 2\omega_2$ has partition form $(2,2,0)$.
The minuscule weight $\omega_1$ has $W_0$-orbit
$W_0 \omega_1 = \{e_1, e_2, e_3\}$ (in $\Z^3$).

\subsection{vDEZ coefficient formula (2.11)}

For simply-laced $A_{n-1}$ with a single parameter $t$, vDEZ Eq.~(2.11) reads
\[
V_{\lambda, \nu}(t) = \prod_{\substack{\beta \in R_0^+\\ \langle\lambda,\beta\rangle=0\\ \langle\nu,\beta\rangle>0}} \frac{1 - t\, \hat e_t(\beta)}{1 - \hat e_t(\beta)}
\;\cdot\;
\prod_{\substack{\beta \in R_0^+\\ \langle\lambda,\beta\rangle=c\\ \langle\nu,\beta\rangle<0}} \frac{1 - t\, \hat h_t\, \hat e_t(-\beta)}{1 - \hat h_t\, \hat e_t(-\beta)},
\]
where $\hat e_t(\eta) = \prod_{\beta \in R_0^+} t^{\langle \eta, \beta\rangle/2}$
and $\hat h_t = t \cdot \hat e_t(-\alpha_0) = t \cdot \hat e_t(\theta)$.

Direct calculation: in $A_2$, $\hat e_t(\alpha_1) = t$, $\hat e_t(\alpha_2) = t$,
$\hat e_t(\theta) = t^2$, hence $\hat h_t = t \cdot t^2 = t^3$.

\subsection{Explicit computation}

For $\lambda = (2, 2, 0)$ and $\nu \in W_0\omega_1$, we tabulate $\langle \lambda, \beta\rangle$,
$\langle \nu, \beta\rangle$, and $V_{\lambda,\nu}(t)$:

\begin{center}
\begin{tabular}{|c|c|c|c|c|c|c|}
\hline
$\nu$ & $\lambda + \nu$ & in $P_c$? & $\langle\lambda,\alpha_1\rangle,\langle\nu,\alpha_1\rangle$ & $\langle\lambda,\alpha_2\rangle,\langle\nu,\alpha_2\rangle$ & $\langle\lambda,\theta\rangle,\langle\nu,\theta\rangle$ & $V_{\lambda,\nu}(t)$ \\
\hline
$e_1$ & $(3,2,0)$ & YES     & $0, +1$ & $2, 0$ & $2, +1$ & $\frac{1-t \cdot t}{1-t} = 1 + t$ \\
$e_2$ & $(2,3,0)$ & no (nondom) & $0, -1$ & $2, +1$ & $2, 0$ & --- \\
$e_3$ & $(2,2,1)$ & YES     & $0, 0$ & $2, -1$ & $2, -1$ & $1$ (empty product) \\
\hline
\end{tabular}
\end{center}

Explanation of row 1: the only $\beta$ satisfying $\langle\lambda,\beta\rangle = 0$ is
$\beta = \alpha_1$. It has $\langle\nu = e_1, \alpha_1\rangle = 1 > 0$, so
contributes $(1 - t\, \hat e_t(\alpha_1))/(1 - \hat e_t(\alpha_1))
= (1 - t^2)/(1-t) = 1+t$. No $\beta$ satisfies $\langle\lambda,\beta\rangle = c = 3$
(the values are $0, 2, 2$), so the second product is empty.

Row 3: no $\beta$ satisfies BOTH $\langle\lambda,\beta\rangle = 0$ and
$\langle\nu = e_3, \beta\rangle > 0$ (only $\alpha_1$ passes the first, but
$\langle e_3, \alpha_1\rangle = 0$). Second product is empty for the same
level-alcove reason.

\subsection{Resulting Pieri identity}

Substituting into vDEZ Theorem~2.4 (with $U_{\lambda,\omega_1} = 0$ since $\omega_1$ is minuscule):

\begin{proposition}[Explicit Pieri identity for cylindric HL / periodic spherical, $A_2$ level 3]
\label{prop:explicit-pieri}
In the periodic Macdonald spherical function basis of vDEZ:
\[
m_{\omega_1}\cdot M^{(c=3)}_{2\omega_2}
\;=\; (1+t)\cdot M^{(c=3)}_{\omega_1 + 2\omega_2}
\;+\; M^{(c=3)}_{\omega_2}.
\]
Equivalently, in the ordinary partition notation:
\[
m_{\omega_1}\cdot M^{(c=3)}_{(2,2,0)}
\;=\; (1+t)\cdot M^{(c=3)}_{(3,2,0)}
\;+\; M^{(c=3)}_{(2,2,1)}.
\]
\end{proposition}

\begin{remark}[Cylindric wrap-around]
The identification $(2,2,1) \leftrightarrow \omega_2$ (rather than
$\omega_1 + \omega_2$) reflects a translation by the sum-preserving element
$(1,1,1)$: as an $A_2$ weight, $(2,2,1) \equiv (1,1,0) = \omega_2$. This is
the ``cylindric wrap-around'': adding the minuscule box to $(2,2,0)$ can
either land at $(3,2,0)$ (dominant, on the boundary $\langle\lambda,\theta\rangle=3=c$
with the extra factor $1+t$) or wrap around to $(2,2,1) \equiv \omega_2$
(dominant, deep in the alcove, coefficient $1$).
\end{remark}

\begin{remark}[Linear limit]
Taking $c \to \infty$ in this Pieri identity would require $M^{(c)}_\lambda \to P_\lambda$
under appropriate normalisation. The first-order coefficient at the ``boundary''
term $(3,2,0)$ decouples from $c$ in vDEZ Eq.~(2.11) only after this
normalisation. In our ordinary-HL computation
(Section~\ref{ssec:hl-pieri-comparison}), we verified
\[
e_1 P_{(2,2,0)} \;=\; P_{(3,2,0)} + P_{(2,2,1)}
\]
in the ordinary HL polynomial ring, with coefficient $1$ (not $1+t$) at
$P_{(3,2,0)}$. So the extra $t$ in Proposition~\ref{prop:explicit-pieri} is
a genuine cylindric level-$c$ effect.
\end{remark}

\section{Independent cross-check: Warnaar Theorem~1.1 (unconditional)}
\label{sec:warnaar-crosscheck}

For $\mu = (k^r)$ rectangular, Warnaar \cite{Warnaar2026} proved unconditionally
the affine dual Jacobi--Trudi identity
\[
P_{(k^r)}(x;t) \;=\; \sum_{y \in Q} \det_{1 \le i, j \le k}\bigl(t^{k\binom{y_i}{2} + i y_i}\, e_{r - i + j - k y_i}(x)\bigr),
\]
where $Q = \{y \in \Z^k : \sum y_i = 0\}$ is the root lattice of the affine
Weyl group of type $\wt A_{k-1}$. The LHS is the \emph{ordinary} $P_{(k^r)}$
(not cylindric), evaluated in $x = (x_1, \ldots, x_n)$; the RHS is the
affine-lattice sum.

For our test case $\mu = (2, 2) = (k^r)$ with $k = 2, r = 2$: Warnaar's identity
provides a determinantal formula for $P_{(2,2)}$ that could be cross-checked
against my HL engine. We defer full numerical cross-check to a follow-up
cycle. The key structural relevance: Warnaar's identity is an
\emph{unconditional affine-Weyl-lattice-sum} expansion of a specific
$P_\mu$; a positive answer to the sharper conjecture $(\dgg)$ would
subsume it (up to the identification of coefficients).

\section{The remaining gap}
\label{sec:gap}

The proof of Theorem~\ref{thm:pieri-target} is complete modulo two items:

\begin{enumerate}[label=(\alph*)]
\item \emph{Conjecture~\ref{conj:daggger} $(\dgg)$ itself.} This is the affine
Cherednik--Ram identity at the trivial-Pieri level: the level-$k$ analogue of
Proposition~\ref{prop:linear-cr}. Its linear limit ($k \to \infty$) reduces
to Proposition~\ref{prop:linear-cr}, which is proved. To directly verify
$(\dgg)$ numerically requires an implementation of the extended affine
Demazure--Lusztig operators $\wt T_j$ on the cylindric HL polynomial ring
at a specific level $k$ — not yet built.

\item \emph{The identification of cylindric HL $P^{(c,k)}_\mu$ with periodic Macdonald
spherical functions $M^{(c)}_\mu$.} van~Diejen--Emsiz--Zurri\'an
\cite[Remark~4.3 sequel and preceding paragraph]{vDEZ2021} note that Corollary~4.2
``reproduces the affine Pieri rule for cylindric Hall--Littlewood polynomials
due to Korff'', but the identification $P^{(c,k)}_\mu = c_\mu \cdot M^{(c)}_\mu$
requires tracing through Korff \cite{Korff2013} carefully. Once this is
in hand, vDEZ Cor.~4.2 provides the Pieri identity for the RHS of the target
theorem; whether that identity \emph{proves} $(\dgg)$ requires additional
work (Cor.~4.2 is a Pieri identity, not a symmetrizer identity).
\end{enumerate}

Neither gap is addressed here. The value of Theorem~\ref{thm:pieri-target}
is that it \emph{reduces} the Pieri-level affine CR to the ``trivial'' affine CR
$(\dgg)$ via a one-line centrality argument, and thus \emph{isolates} the
non-trivial content: the trivial-Pieri identity itself.

\section{Ship notes}

Companion Python scripts:
\begin{itemize}[nosep]
\item \verb|verify_linear_pieri.py| — 14/14 pass (Section~\ref{ssec:linear-limit-verif}).
\item \verb|verify_affine_commutation.py| — 95/95 pass at $q=1$, 0/20 pass at generic $q$
  (Section~\ref{ssec:t0-centrality}).
\item \verb|verify_vdez_specialisation.py| — reproduces
  Proposition~\ref{prop:explicit-pieri} from scratch.
\end{itemize}

All scripts sit at \verb|~/projects/scratch/2026-07-24-pieri-affine-CR/|.
Builds require \verb|demazure_engine.py|, \verb|hecke_HL.py|,
\verb|hall_littlewood_engine.py| from \verb|~/projects/scratch/| (2026-07-20 vintage).

\begin{thebibliography}{9}
\bibitem{Cherednik}
I.~Cherednik, \emph{Double affine Hecke algebras}, LMS Lecture Note Series 319,
Cambridge Univ.\ Press, 2005.

\bibitem{ClioLinear}
Clio (2026-07-20), \emph{Demazure operators, Hall--Littlewood via Hecke, and the transfer-op vs Demazure comparison}, \verb|~/projects/proofs/2026-07-20-demazure-basics-and-cherednik-ram-verification.pdf|.

\bibitem{CR}
I.~Cherednik and A.~Ram, \emph{Double affine Hecke algebras and Macdonald's conjectures}, Ann.\ Math.\ 141 (1995), 191--216 (linear CR appears as the $q \to 0$ specialisation of Cherednik's Macdonald construction; cf.\ \cite{Macdonald2003}).

\bibitem{Korff2013}
C.~Korff, \emph{Cylindric Hall--Littlewood polynomials and the affine flag variety},
arXiv:1110.6356, Comm.\ Math.\ Phys.\ 318 (2013), 173--246.

\bibitem{Macdonald2003}
I.~G.~Macdonald, \emph{Affine Hecke Algebras and Orthogonal Polynomials},
Cambridge Tracts in Math.\ 157, Cambridge Univ.\ Press, 2003.

\bibitem{vDEZ2021}
J.F.~van~Diejen, E.~Emsiz, I.N.~Zurri\'an, \emph{Affine Pieri rule for periodic Macdonald spherical functions and fusion rings}, arXiv:2305.01931, Adv.\ Math.\ 392 (2021), 108027.

\bibitem{Warnaar2026}
S.O.~Warnaar, \emph{Affine Jacobi--Trudi identities and $q, t$-Rogers--Ramanujan identities}, arXiv:2511.17034, SIGMA 22 (2026), 062.

\end{thebibliography}

\end{document}
